\documentclass[final]{beamer}
\usepackage[size=a0,orientation=portrait,scale=1.08]{beamerposter}
\usepackage{graphicx,booktabs,amsmath,ragged2e}
\setbeamertemplate{navigation symbols}{}
\definecolor{Accent}{RGB}{158,31,36}
\definecolor{Ink}{RGB}{34,34,34}
\definecolor{Rule}{RGB}{224,224,224}
\setbeamercolor{background canvas}{bg=white}
\setbeamercolor{normal text}{fg=Ink}
\title{Academic Poster Preview}
\subtitle{Minimal portrait layout with large quantitative callouts}
\author{Author Names}
\institute{Institution / Department}
\newcommand{\sectionlabel}[1]{\vspace{.45cm}{\Large\bfseries\color{Accent}#1}\par\vspace{.25cm}}
\newcommand{\metric}[2]{\centering{\fontsize{64}{72}\selectfont\bfseries\color{Accent}#1}\par{\Large\color{gray}#2}}
\begin{document}
\begin{frame}[t]
\vspace{1.1cm}
\begin{columns}[c,totalwidth=.88\paperwidth]
\begin{column}{.74\paperwidth}
{\VeryHuge\bfseries\color{Ink}\inserttitle\par}
\vspace{.35cm}
{\LARGE\color{gray}\insertsubtitle\par}
\vspace{.2cm}
{\Large\color{gray}\insertauthor\quad|\quad\insertinstitute}
\end{column}
\begin{column}{.12\paperwidth}
\centering\includegraphics[width=.85\linewidth]{logo-or-qr}
\end{column}
\end{columns}
\vspace{.55cm}
\noindent\hspace{.06\paperwidth}\textcolor{Accent}{\rule{.88\paperwidth}{3pt}}
\vspace{.9cm}
\begin{columns}[c,totalwidth=.88\paperwidth]
\begin{column}{.24\paperwidth}\metric{--\%}{primary outcome}\end{column}
\begin{column}{.24\paperwidth}\metric{n=--}{samples or episodes}\end{column}
\begin{column}{.24\paperwidth}\metric{--x}{improvement}\end{column}
\begin{column}{.24\paperwidth}\metric{$p<--$}{statistical signal}\end{column}
\end{columns}
\vspace{.7cm}
\noindent\hspace{.06\paperwidth}\textcolor{Rule}{\rule{.88\paperwidth}{1pt}}
\vspace{.7cm}
\begin{columns}[t,totalwidth=.88\paperwidth]
\begin{column}{.42\paperwidth}
\sectionlabel{Motivation and Goal}
\justifying Concise motivation, research gap, and hypothesis. Keep prose short and use whitespace to emphasize the central message.
\sectionlabel{Theory and Formula}
Key equation space: $y=f(x;\theta)$, assumptions, symbol definitions, and how measured variables connect to the model.
\sectionlabel{Model and Method}
\centering\includegraphics[width=.9\linewidth,height=.22\paperheight,keepaspectratio]{workflow-image}
\end{column}
\hfill
\begin{column}{.42\paperwidth}
\sectionlabel{Experimental Workflow}
\centering\includegraphics[width=.9\linewidth,height=.2\paperheight,keepaspectratio]{schematic-image}
\sectionlabel{Quantitative Results}
\centering\includegraphics[width=.9\linewidth,height=.18\paperheight,keepaspectratio]{result-panel}
\sectionlabel{Interpretation}
\justifying Compact takeaways, limitations, future work, references, contact details, and approved QR/logo assets.
\end{column}
\end{columns}
\end{frame}
\end{document}
